% GENERATED FILE. Edit canonical lessons, then run npm run generate:books.
% canonical-source-hash: fnv1a64:ca94d5f8843f0795
% canonical-lessons: PT-C431-radiador, PT-C431-aquecedor, PT-C431-esquentador, PT-C431-termostato, PT-C431-arejar

\chapter{Calor e frio no quarto}
\label{ch:pt-calor-e-frio-no-quarto}
\hlchaptermodality{431}

\begin{chapteropening}
\textbf{By the end of this chapter:} \emph{I can complain about the temperature: the radiator, the heater, the water heater, the thermostat and airing the room.}

Complete the last lesson of chapter 431: I can complain about the temperature: the radiator, the heater, the water heater, the thermostat and airing the room.
\end{chapteropening}

\section[o radiador]{o radiador — the radiator}

\label{lesson:PT-C431-radiador}



\begin{quote}
Before the new one: say the Portuguese for to unblock, then the Portuguese for the toilet cistern.
\end{quote}

\subsection*{You'll want to know: o radiador}
\textbf{o radiador} — \textquotedblleft{}the radiator\textquotedblright{}.

\begin{quote}
O radiador do quarto está frio.
\end{quote}

\begin{quote}
\emph{The radiator in the room is cold.}
\end{quote}

A car has \emph{um radiador} too.

\begin{cousinweb}
\emph{radiador} is built on \emph{radiar} 'to give off rays', a learned borrowing from Latin \emph{radiāre}, on \emph{radius} 'ray, spoke' (Lewis \& Short).
\end{cousinweb}

\begin{grammarlens}[title={What you've built}]
One more word for telling, explaining and describing at B1.
\end{grammarlens}

\subsection*{Your turn}
\begin{itemize}
\raggedright
  \item \emph{Say it:} \emph{o radiador}
  \item \emph{Say it:} \emph{o radiador}, once more
  \item \emph{Say it:} a sentence of your own with \emph{o radiador}
  \item \emph{Recall:} say the Portuguese for to unblock, then the Portuguese for the toilet cistern, then say \emph{o radiador} again
\end{itemize}

\begin{tcolorbox}[breakable,colback=teal!4,colframe=teal!35!black,title={Before you move on}]
What does o radiador mean? (\textquotedblleft{}The radiator\textquotedblright{}.) Say it once more.
\end{tcolorbox}
\section[o aquecedor]{o aquecedor — the heater}

\label{lesson:PT-C431-aquecedor}



\begin{quote}
Before the new one: say the Portuguese for the toilet cistern, then the Portuguese for the radiator.
\end{quote}

\subsection*{You'll want to know: o aquecedor}
\textbf{o aquecedor} — \textquotedblleft{}the heater\textquotedblright{}.

\begin{quote}
Pedimos um aquecedor porque o quarto estava frio.
\end{quote}

\begin{quote}
\emph{We asked for a heater because the room was cold.}
\end{quote}

\emph{Um aquecedor elétrico} — an electric heater.

\begin{cousinweb}
\emph{aquecedor} is built on \emph{aquecer} 'to warm', which goes back to Latin \emph{calēscere} 'to grow warm', on \emph{calēre} 'to be warm' (Lewis \& Short).
\end{cousinweb}

\begin{grammarlens}[title={What you've built}]
One more word for telling, explaining and describing at B1.
\end{grammarlens}

\subsection*{Your turn}
\begin{itemize}
\raggedright
  \item \emph{Say it:} \emph{o aquecedor}
  \item \emph{Say it:} \emph{o aquecedor}, once more
  \item \emph{Say it:} a sentence of your own with \emph{o aquecedor}
  \item \emph{Recall:} say the Portuguese for the toilet cistern, then the Portuguese for the radiator, then say \emph{o aquecedor} again
\end{itemize}

\begin{tcolorbox}[breakable,colback=teal!4,colframe=teal!35!black,title={Before you move on}]
What does o aquecedor mean? (\textquotedblleft{}The heater\textquotedblright{}.) Say it once more.
\end{tcolorbox}
\section[o esquentador]{o esquentador — the water heater (gas)}

\label{lesson:PT-C431-esquentador}



\begin{quote}
Before the new one: say the Portuguese for the radiator, then the Portuguese for the heater.
\end{quote}

\subsection*{You'll want to know: o esquentador}
\textbf{o esquentador} — \textquotedblleft{}the water heater (gas)\textquotedblright{}.

\begin{quote}
O esquentador avariou e não há água quente.
\end{quote}

\begin{quote}
\emph{The water heater broke down and there is no hot water.}
\end{quote}

In Brazil \emph{o aquecedor} covers both. \emph{Esquentar} — to heat up.

\begin{cousinweb}
\emph{esquentador} is built on \emph{esquentar} 'to heat', from \emph{es-} and \emph{quente} 'hot' (EtymDB); \emph{quente} is inherited from Latin \emph{calēns} 'warm'.
\end{cousinweb}

\begin{grammarlens}[title={What you've built}]
One more word for telling, explaining and describing at B1.
\end{grammarlens}

\subsection*{Your turn}
\begin{itemize}
\raggedright
  \item \emph{Say it:} \emph{o esquentador}
  \item \emph{Say it:} \emph{o esquentador}, once more
  \item \emph{Say it:} a sentence of your own with \emph{o esquentador}
  \item \emph{Recall:} say the Portuguese for the radiator, then the Portuguese for the heater, then say \emph{o esquentador} again
\end{itemize}

\begin{tcolorbox}[breakable,colback=teal!4,colframe=teal!35!black,title={Before you move on}]
What does o esquentador mean? (\textquotedblleft{}The water heater (gas)\textquotedblright{}.) Say it once more.
\end{tcolorbox}
\section[o termóstato]{o termóstato — the thermostat}

\label{lesson:PT-C431-termostato}



\begin{quote}
Before the new one: say the Portuguese for the heater, then the Portuguese for the water heater.
\end{quote}

\subsection*{You'll want to know: o termóstato}
\textbf{o termóstato} — \textquotedblleft{}the thermostat\textquotedblright{}.

\begin{quote}
Baixei o termóstato para vinte graus.
\end{quote}

\begin{quote}
\emph{I turned the thermostat down to twenty degrees.}
\end{quote}

In Brazil: \emph{o termostato}.

\begin{cousinweb}
\emph{termóstato} is a learned coinage from Greek \emph{thermós} 'hot' and \emph{statós} 'standing still' (Lewis \& Short has Latin \emph{thermae} 'hot baths', from the same Greek word): it keeps the heat steady.
\end{cousinweb}

\begin{grammarlens}[title={What you've built}]
One more word for telling, explaining and describing at B1.
\end{grammarlens}

\subsection*{Your turn}
\begin{itemize}
\raggedright
  \item \emph{Say it:} \emph{o termóstato}
  \item \emph{Say it:} \emph{o termóstato}, once more
  \item \emph{Say it:} a sentence of your own with \emph{o termóstato}
  \item \emph{Recall:} say the Portuguese for the heater, then the Portuguese for the water heater, then say \emph{o termóstato} again
\end{itemize}

\begin{tcolorbox}[breakable,colback=teal!4,colframe=teal!35!black,title={Before you move on}]
What does o termóstato mean? (\textquotedblleft{}The thermostat\textquotedblright{}.) Say it once more.
\end{tcolorbox}
\section[arejar]{arejar — to air (a room); to get some fresh air}

\label{lesson:PT-C431-arejar}



\begin{quote}
Before the new one: say the Portuguese for the water heater, then the Portuguese for the thermostat.
\end{quote}

\subsection*{You'll want to know: arejar}
\textbf{arejar} — \textquotedblleft{}to air (a room); to get some fresh air\textquotedblright{}.

\begin{quote}
Abro a janela para arejar o quarto.
\end{quote}

\begin{quote}
\emph{I open the window to air the room.}
\end{quote}

\emph{Arejar a cabeça} — to clear your head.

\begin{cousinweb}
\emph{arejar} is built on \emph{ar} 'air', which is inherited from Latin \emph{āēr}, from Greek \emph{a\'{\={e}}r} (Lewis \& Short).
\end{cousinweb}

\begin{grammarlens}[title={What you've built}]
One more word for telling, explaining and describing at B1.
\end{grammarlens}

\subsection*{Your turn}
\begin{itemize}
\raggedright
  \item \emph{Say it:} \emph{arejar}
  \item \emph{Say it:} \emph{arejar}, once more
  \item \emph{Say it:} a sentence of your own with \emph{arejar}
  \item \emph{Recall:} say the Portuguese for the water heater, then the Portuguese for the thermostat, then say \emph{arejar} again
\end{itemize}

\begin{tcolorbox}[breakable,colback=teal!4,colframe=teal!35!black,title={Before you move on}]
What does arejar mean? (\textquotedblleft{}To air (a room); to get some fresh air\textquotedblright{}.) Say it once more.
\end{tcolorbox}
